\documentclass[../../main.tex]{subfiles}
\begin{document}
\begin{oframed}
  \subfile{problem}
\end{oframed}

{\bf [解答]}

$n$の時の$a,b$を$a_n,b_n$と書くと，適当な$n-2$次式$P_n(x)$を用いて
\begin{align}
  x^n=P_{n}(x)(x^2-2x-1)+a_nx+b_n \label{eq:1}
\end{align}
と書ける．両辺に$x$をかけて
\begin{align}
  x^{n+1}=\left[xP_{n}(x)+a_n\right](x^2-2x-1)+(2a_n+b_n)x+a_n
\end{align}
だから，
\begin{align}
  \begin{cases}
    a_{n+1} = 2a_n+b_n \\
    b_{n+1} = a_n
  \end{cases} \label{eq:2}
\end{align}
となる．初項は
\begin{align}
  a_1 = 1, b_1 = 0 \label{eq:3}
\end{align}

以下$a_n,b_n\in\mathbb{Z}$かつ$a_n$と$b_n$が互いに素であることを帰納的に示す．
$n=1$の時は\cref{eq:3}より成立するから，以下$n=k$での成立を仮定する．$n=k+1$の時は\cref{eq:2}より
\begin{align}
  a_{k+1}=2a_k+b_k\in\mathbb{Z}，\qquad b_{k+1}=a_k \in\mathbb{Z} \label{eq:4}
\end{align}
だから$a_{k+1},b_{k+1}$はともに整数である．
次に$a_{k+1}$，$b_{k+1}$が共通の素因数$d$を持つとすると，\cref{eq:4}より$a_k$, 従って$b_k$も$d$で割り切れ仮定に反する．よって$a_{k+1}$, $b_{k+1}$は互いに素である．

以上から$n=k+1$でも成立し，数学的帰納法により題意は示された．$\blacksquare$

\newpage

\end{document}
